\documentclass[12pt]{article}
\usepackage{geometry}
\geometry{a4paper, portrait, margin=1in}

% --- Math Packages ---
\usepackage{amsmath}
\usepackage{amssymb}
\usepackage{amsfonts}
\usepackage{amsthm}

% --- Graphics & Diagrams ---
\usepackage{graphicx}
\usepackage{tikz}
\usetikzlibrary{shapes, arrows, positioning}

% --- Formatting ---
\usepackage{parskip}
\usepackage[hidelinks]{hyperref}
\usepackage{natbib}
\usepackage{enumitem}
\pagestyle{plain}

\title{MPhil Thesis: Introduction}
\author{Mohammad-Shah Karim}
\date{April 2026}

\begin{document}

\maketitle

\section{Introduction}

Pharmaceutical innovation is central towards long-run improvements in health and welfare. New drugs can improve quantity of life while changing the treatment possibilities which are available for physicians to administer on patients. At the same time, pharmaceutical innovation is unusually sensitive to policy. Development of a new drug requires a large upfront investment with a long development horizon including patent applications and clinical trials, with only a high probability of failure. Since the returns to a successful drug accrue over multiple years and depend on both prices and the volumes that prevail within the market, the incentives to undertake the investment to develop a drug are shaped not only by opportunity, but also by the economic and institutional environment present which determine the expected future profits once the drug reaches the market.

This makes the pharmaceutical market within the United States of America an ideal setting for studying how public policy shapes scientific opportunity. Federal and state governments are now responsible for roughly half of all prescription drug spending (INSERT SOURCE HERE), both as direct purchasers through programmes such as Medicaid and Medicare, and also indirectly through the tax preferences and regulatory rules that govern private insurance. Medicaid is especially important in this respect - it is a large public buyer, but it does not operate as a standard private insurer. Manufacturers that wish to have Medicaid coverage must participate in the Medicaid Drug rebate Program, and states can use Preferred Drug Lists (PDLs) in order to negotiate additional rebates in exchange for preferred formulary placement. Given this, Medicaid demand passes through a public procurement system with statutory and supplement rebates in addition to state level formulary bargaining, thus it cannot be treated as additional market demand.

This thesis examines the consequences of one of such reforms for pharmaceutical innovation - the 2014 Medicaid expansion under the the Patient Protection and Affordable Care Act (ACA). The reform produced a well-documented expansion in the size of the Medicaid market, with population coverage increasing by close to twenty million enrollees int he years following the policy change (INSERT SOURCE HERE). Under standard market-sized frameworks in the pharmaceutical sector this should increase innovation incentives for drug-classes more exposed to Medicaid demand \citep{finkelstein2004static, dubois2015market}. If firms expect to sell more units after a successful round of innovation, the expected returns to R\&D rise. However, the expansion of Medicaid took place inside a procurement framework that combined two features of that are unusual under the standards of a private insurer. A statutory rebate schedule that mechanically transfers a share of any list-price revenue back to the state, and a state-administered formulary system in which Preferred Drug Lists (PDLs) determine which products are reimbursed without prior authorisation. Thus the institutional features meant that an expansion of Medicaid was not simply a positive demand shock for pharmaceutical firm.

The Medicaid setting therefore complicates the prediction under standard pharmaceutical market-sized frameworks. The ACA did not only expand the number of beneficiaries under Medicaid, but it also raised the statutory rebate percentages and extended rebate obligations to Medicaid managed care. Therefore, the reform increased the size of the public market while also increasing the share of revenue back to the state. Additionally, as Medicaid demand grows, preferred placement on a state PDL becomes more valuable. States can use this to negotiate deeper rebates from manufacturers. The same policy may therefore generate two opposing effects. On the one hand, Medicaid expansion raises demand. On the other hand, it may increase pricing pressure and reduce the surplus that firms can capture from that demand.

This thesis therefore raises the question of whether the ACA-induced Medicaid expansion raised pharmaceutical innovation, and whether the response dependent on a firm's exposure to the Medicaid procurement institutions. Put simply, was the expected positive relationship between market size and innovation weakened when the additional demand flows through a state-wide monopsony. Many current drug-pricing reforms, including recent Medicare negotiation policies, rely on the same basic trade-off: lowering public spending and consumer prices may increase static welfare, but it may also affect the expected returns to future innovation, raising the relevance of this question


\newpage
\bibliographystyle{apalike}
\bibliography{references}

\end{document}
